\section{Some Concepts}
\begin{definition}[Inverse Function]
	Let $f : \mathcal{X} \to \mathcal{Y}$ be an uncertain function on the arithmetic space $\alpha$ with $\alpha \in \set{1,m,m'}$, $X \in \mathcal{X}$. We define the inverse operation of $f$ on the arithmetic space $\alpha$ as:
	\[
		(f^{-1}(X))_{\alpha} := \s{A \in \tsbd{K}: (f(A))_\alpha = X}
	\]
\end{definition}
\begin{example}
	\begin{enumerate}
		\item Given $f(X)=\kgdd{X^2}$, we define: \[\kgdd{\sqrt{X}} = \s{A \in \tsbd{K} : \kgdd{X^2} = X}\]
		\item Given $f(X)=\kgdd{\sin{X}}$, we define: \[\kgdd{\arcsin{X}} = \s{A \in \tsbd{K} : \kgdd{sin(A)} = X}\]
	\end{enumerate}
\end{example}
\begin{example}
	Set $N_0 = \s{1,2,5}$ and $f(X)=X^{N_0}$, we define: \[\sqrt[N_0]{X} = \s{A \in \tsbd{K} : A^{N_0} = X}\]
For instance \[\sqrt[N_0]{\s{1, 2, 3, 4, 9, 32, 243}} = \s{1,2,3}\]
\end{example}
\begin{example}
	Given $f(X)=\kgdtmr{X^2}$, we define: \[\kgdtmr{\sqrt{X}} = \s{A \in \tsbd{K} : \kgdtmr{A^2} = X}\]
	For instance \[\kgdtmr{\sqrt{\s{1, 2, 4}}} = \s{\s{1,2},\s{-1,-2}}\]
\end{example}
\begin{definition}[Graph of an Uncertain Function at a Point]
	Let $f : \mathcal{X} \to \mathcal{Y}$ be an uncertain function, $A \in \mathcal{X}$, we define the graph of the uncertain function $f$ at point $A$ as:
	\[pgr_f(A):= A \times f(A)\]
\end{definition}
\begin{definition}[Graph of an Uncertain Function]
	Let $\mathcal{X} ,\mathcal{Y}$ be subsets of $\tsbd{K}$ where $\mathbb{K}$ is an arbitrary field.
	Let $f : \mathcal{X} \to \mathcal{Y}$ be an uncertain function, $A \in \mathcal{X}$, we define the graph of the uncertain function $f$ as:
	\[\gr{f}:= \set{\pgr_f(X) : X \in \mathcal{X}}\]
\end{definition}
\begin{example}
	For $f = X+X$, we have \[\pgr_f{\s{1,2}} = \s{(1,2),(1,3),(1,4),(2,2),(2,3),(2,4)}\]
\end{example}
\begin{example}
	Given $$X = \dfrac{1}{10}\s{10..20} $$ and the function $$m(X) = X^{\s{2,3}}+5X $$ we obtain the result that $m(X)$ equals
	$\{$6, 6.21, 6.331, 6.44, 6.5, 6.69, 6.71, 6.728, 6.831, 6.94, 6.96, 7, 7.19, 7.197, 7.21, 7.228, 7.25, 7.331, 7.44, 7.46, 7.5, 7.56, 7.69, 7.697, 7.71, 7.728, 7.744, 7.75, 7.831, 7.89, 7.94, 7.96, 8, 8.06, 8.19, 8.197, 8.21, 8.228, 8.24, 8.244, 8.25, 8.331, 8.375, 8.39, 8.44, 8.46, 8.5, 8.56, 8.61, 8.69, 8.697, 8.71, 8.728, 8.74, 8.744, 8.75, 8.831, 8.875, 8.89, 8.94, 8.96, 9, 9.06, 9.096, 9.11, 9.19, 9.197, 9.21, 9.228, 9.24, 9.244, 9.25, 9.331, 9.375, 9.39, 9.44, 9.46, 9.5, 9.56, 9.596, 9.61, 9.69, 9.697, 9.71, 9.728, 9.74, 9.744, 9.75, 9.831, 9.875, 9.89, 9.913, 9.94, 9.96, 1E+1, 10.06, 10.096, 10.11, 10.19, 10.197, 10.21, 10.228, 10.24, 10.244, 10.25, 10.331, 10.375, 10.39, 10.413, 10.44, 10.46, 10.5, 10.56, 10.596, 10.61, 10.69, 10.697, 10.71, 10.728, 10.74, 10.744, 10.75, 10.831, 10.832, 10.875, 10.89, 10.913, 10.94, 10.96, 11, 11.06, 11.096, 11.11, 11.19, 11.197, 11.21, 11.228, 11.24, 11.244, 11.25, 11.331, 11.332, 11.375, 11.39, 11.413, 11.44, 11.46, 11.5, 11.56, 11.596, 11.61, 11.69, 11.697, 11.728, 11.74, 11.744, 11.75, 11.832, 11.859, 11.875, 11.89, 11.913, 11.96, 12, 12.06, 12.096, 12.11, 12.197, 12.24, 12.244, 12.25, 12.332, 12.359, 12.375, 12.39, 12.413, 12.5, 12.56, 12.596, 12.61, 12.74, 12.744, 12.832, 12.859, 12.875, 12.89, 12.913, 13, 13.096, 13.11, 13.24, 13.332, 13.359, 13.375, 13.413, 13.5, 13.596, 13.61, 13.832, 13.859, 13.913, 14, 14.096, 14.332, 14.359, 14.413, 14.5, 14.832, 14.859, 14.913, 15, 15.332, 15.359, 15.5, 15.832, 15.859, 16, 16.359, 16.5, 16.859, 17, 17.5, 18$\}_u$
	\newline \noindent Number of elements: 223, and this set exhibits a structure that we see as similar to the interval $[6,18]$.
	\newline \noindent If we take $X = \kbd{[1,2]}$, then we have $m(X) = \kbd{[6,18]}$.
\end{example}
\begin{example}
	Given $X \in \tsbd{R}$ where all elements are positive integers, we have:
	$$X^3 \mod 3 = \s{0,1,2}$$
	$$X^3 \mod 3 = \s{0,1,3}$$
	$$X^3 \mod \s{3,4} = \s{0,1,2,3}$$
	These are constant uncertain functions.
\end{example}
\begin{definition}[Trace]
	Let $f$ be an uncertain function from $\mathcal{X}$ into $\mathcal{Y}$ mapping $X \mapsto Y = f(X)$.
	\newline \noindent
	Considering $y \in Y = f(X)$, we define the trace of $y$ in the expression of $f(X)$, denoted as $\Tr(y \in f(X))$, as the set of all formal expressions whose evaluated result equals $y$, obtained by computing $f(X)$.
\end{definition}
\begin{definition}[Uncertain Equation]
	Let $f : \mathcal{X}\to \mathcal{Y}$ be an uncertain function, $C$ be an uncertain constant; an uncertain equation with unknown $X$ is an equation of the form: \[f(X) = C\]
\end{definition}
\begin{example}
	Given the uncertain equation \[\kgdd{X^2+5X+6} = 0 \]
	It is easy to see that $X = \s{-2,-3}$ is a solution to this equation. \\
	However, $X = \s{-2,-3}$ is not a solution to the equation:
	\[X^2+5X+6 = 0 \]
	The equation $X^2+5X+6 = 0$ has solutions $X=-2$ and $X=-3$.
\end{example}
\begin{definition}[Weak Solution of an Uncertain Equation]
	Let $f : \mathcal{X} \to \mathcal{Y}$ be an uncertain function, $C \in \tsbd{K}$. 
	A solution $X_0 \in \tsbd{K}$ is called a weak solution of the uncertain equation $f(X) = C$ if $f(X_0)$ is defined and $f(X_0) \subsetneq C$. If $f(X_0)$ and $C$ are finite-element uncertain numbers, we set: \[\gamma = \frac{\card(f(X_0))}{\card(C)} \in (0,1),\] and we write $(f(X_0) = C)_\gamma$, where the equality relation here is a weak equality relation. For example, if $\gamma = 0.5$, we write $\gtcl{f(X_0) = C}{0.5}$.
\end{definition}
\begin{theorem}
	Let $f : \mathcal{X} \to \mathcal{Y}$ be an uncertain function, $C \in \tsbd{K}$. If the uncertain equation $f(X) = C$ has a solution, then there exist $A,B \in \tsbd{K}$ such that $\gtcl{f(A) = C}{\gamma}$ and $B$ is a solution of $f(X)=C$ with $A \subseteq B$.
\end{theorem}
\begin{theorem}
	\label{dl:3.2}
	Let $X \in \mathcal{X}$ be an uncertain number with canonical form $X \sim (f_X, d_X)$ and $f: \mathcal{X} \rightarrow \mathcal{Y}$ be a single-variable uncertain mapping. 
	
	Then, the image $Y = f(X) \in \mathcal{Y}$ possesses a canonical form $Y \sim (f_Y, d_Y)$ determined in each arithmetic space as follows:
	
	\begin{enumerate}
		\item \textbf{On the Single-point Space $(\circ)_1$:} 
		Since internal point identity is preserved, the scenario domain of $Y$ remains unchanged:
		\[
		d_Y = d_X
		\]
		The generating function $f_Y: d_Y \rightarrow \mathcal{Y}$ is defined via composition:
		\[
		\forall t \in d_Y, \quad f_Y(t) = f(f_X(t))
		\]
		
		\item \textbf{On the Minkowski Space $(\circ)_m$:} 
		If the expression of the function $f(X)$ contains $k$ occurrences of the variable $X$ ($k \in \mathbb{N}^*$), since each occurrence acts as an independent noise source, the scenario domain expands into a $k$-fold Cartesian product:
		\[
		d_Y = d_X^k = \underbrace{d_X \times d_X \times \dots \times d_X}_{k \text{ times}}
		\]
		The generating function $f_Y: d_Y \rightarrow \mathcal{Y}$ is defined by:
		\[
		\forall (t_1, t_2, \dots, t_k) \in d_Y, \quad f_Y(t_1, t_2, \dots, t_k) = f\left(f_X(t_1), f_X(t_2), \dots, f_X(t_k)\right)
		\]
	\end{enumerate}
\end{theorem}

\begin{remark}
	Theorem \ref{dl:3.2} clarifies the essential boundary between the two arithmetic spaces: the concept of a ``single-variable function'' $f(X)$ in the single-point arithmetic space $(\circ)_1$ preserves the scenario dimension $d_Y = d_X$; whereas in the Minkowski space $(\circ)_m$, the algebraic expression of a single-variable function is deconstructed based on the $k$ occurrences of variable $X$, causing the scenario space to expand into $d_X^k$ corresponding to the Abstract Syntax Tree (AST).
\end{remark}
\begin{theorem}[Solution Characterization of Uncertain Equations]
	Let $f : \mathcal{X} \to \mathcal{Y}$ be an uncertain function and $C \in \mathcal{Y}$ be a target set/constant with canonical representation $C \sim (f_C, d_C)$.
	
	Then, an element $X_0 \in \mathcal{X}$ is a solution of the equation $f(X) = C$ if and only if the value $F = f(X_0) \sim (f_F, d_F)$ satisfies the scenario structure matching condition:
	\[
	\forall t \in d_F, \quad \exists t' \in d_C : \quad f_F(t) = f_C(t')
	\]
\end{theorem}